\minitopic{证言是独立来源吗}我们的大多数知识来自别人：出生日期、遥远城市、科学史以及语言本身。\term{证言}{testimony}不限于法庭陈述，而包括以沟通方式把某内容呈现为真的行为。问题是：听者凭什么从“甲说$p$”到“$p$”？ 还原论要求证言的可信度最终由知觉、记忆和归纳支持，例如说话者过去可靠、当前没有欺骗动机。非还原论则认为，在没有反面迹象时，理解并接受真诚断言本身可给予初步权利；否则儿童和日常交流几乎不可能开始。非还原不等于轻信，因为击败信息仍能撤销默认权利\parencite{burge1993,lackey2008}。 证言知识涉及两方规范。说话者应有适当根据、表达清楚并避免误导；听者应评估能力、诚实、利益冲突和背景一致性。责任不是平均分配：信息不对称越强，专业传播者的说明义务越重。

\minitopic{知识是否能够传递}简单传递模型认为，说话者知道$p$，听者通过可靠证言获得同一知识。但听者有时能从多个各自不知道的报告中整合出知识；说话者有知识时，带有强反证的听者也不能只凭其话获得知识。证言既可传递也可生成认识地位。 拉基的案例还提示：说话者可能没有知识，却从可靠资料背诵真内容，使合理听者获得知识\parencite{lackey2008}。争论取决于我们如何理解群体、媒介和记录在链条中的作用。把证言想成个人头脑之间搬运“知识包裹”，会遗漏制度性保障。

\minitopic{如何识别专家}非专家不能直接以专业内容判断谁更专业，否则已经拥有所求能力。但可以使用第二阶指标：训练与成果、同行认可、解释竞争观点的能力、利益披露、预测记录、对不确定性的校准，以及结论是否经独立方法汇合\parencite{goldman2001}。这些指标不是机械清单；共识可能受制度偏差影响，异议也不因少数身份自动更勇敢或更真。 “诉诸权威”只有在用权威取代本可提供的理由，或权威不相关时才构成谬误。合理信任专家恰恰是有限主体的认识能力之一。成熟的自主不是拒绝依赖，而是能选择、监督和重新分配依赖\parencite{hardwig1985}。

\minitopic{分歧如何影响信心}发现一个与自己证据相当、推理能力相当的\term{认识同侪}{epistemic peer}持相反观点，似乎构成关于自己可能犯错的新证据。调和主义要求降低信心；坚定立场允许在某些条件下维持原判断\parencite{christensen2007,kelly2005}。 “完全同侪”在现实中罕见，但理想化有助于隔离问题。分歧证据与一阶证据可能冲突：我仔细算得答案为17，同侪算得19。若我一律对半调整，会不会忽略自己对计算过程的直接把握？若一律坚持，又会不会把“我看来正确”当成特权？合理回应需考虑独立性、分歧记录、证据共享程度和错误成本。

\begin{comparison}
《论语》既重视“多闻”“择其善者而从之”，也要求学习者反思而非被动接受。师友关系展示一种关系性自主：信任可以培养判断力，而不是必然压制判断力。但古典角色权威也可能固化不平等，因此不能把尊师直接当作现代专家制度的正当化。
\end{comparison}

\minitopic{证言链与来源独立性}五个网站报道同一消息，不一定是五份证据；它们可能都抄自同一匿名帖子。证据数量必须结合来源谱系评价。独立来源在各自可能出错的条件不完全重合时更有增量价值；若共享同一数据库、通讯社或利益冲突，表面汇合会夸大支持。 来源链越长，内容可能越容易失真，但中介也能增加编辑、审查和专业解释。直接目击不天然优于经过核验的机构报告。关键是每一环是否保留出处、是否有纠错机制、最终传播者是否准确表达原来源的确信程度。

\minitopic{不信任也需要理由}批判性思考常被误解为默认拒绝。若听者仅因说话者口音、性别或群体身份降低可信度，这不是自主而是偏见。相反，发现具体利益冲突、专业越界或反复错误记录，构成相关击败者。理性信任和理性不信任都需与证据相称。 阴谋论式怀疑常设置不对称标准：官方证据的任何缺口都被视为造假证明，支持自身的匿名材料却免于审查。纠正办法不是要求信任权威，而是让同一来源标准双向适用，并说明什么证据会改变自己的判断。

\minitopic{案例工作坊：多位专家的共识}四位医生建议相同治疗。若他们独立查看病例并依据不同专长汇合，共识显著增加支持；若后三人只阅读第一人的意见，增量较小；若所有人受同一有缺陷指南约束，一致也可能系统性错误。患者应询问依据是否独立、指南证据等级、替代方案和不确定范围。 专家分歧也不必由非专家“选边”。可以暂停、寻求能解释分歧来源的综合意见，或在可逆性与风险基础上行动。认识成熟表现为把“谁对”转化为“分歧来自数据、价值、模型还是个案适用性”。
